\section{Evaluation Setup}
\label{sec:evaluation}

\paragraph{Protocol.} We use grouped 5-fold cross-validation over repositories: the 53 core
repositories are split into five folds of similar query count, and every fold is tested with
weights learned on the other four, so each repository is scored by a model that never saw it. Only
the 13 weights cross repositories; all signals are always computed within a repository. The grid
search runs inside each fold on training repositories only, and the genetic algorithm is run with
10 seeds per fold. The 8 held-out repositories are never trained on and are scored with every
fold's model.

\paragraph{Metrics.} The primary metric is the \emph{mean reciprocal rank} (MRR), the inverse rank of
the first true answer, which is also the fitness. We also report \emph{mean average precision} (MAP),
which rewards ranking \emph{all} true answers high, and Hit@$k$, the share of queries with a true
answer in the top $k$ (the success rate used by GA-IaCRec~\cite{bessghaier2026iacrec}). Metrics
are averaged per repository and then over repositories, so every repository counts equally. They
are reported for two target sets: \textbf{all targets} (existing blocks and types to add) and the
stricter \textbf{existing blocks only}.

\paragraph{Comparisons.} For RQ1, the GA is compared with a \emph{random ranking} (the
floor for feasibility) and with \emph{random search}, which draws as many random weight vectors as
the genetic algorithm evaluates and keeps the best one; two further checks retrain it with the
labels shuffled within each query, which must destroy what it learns, and with negative weights
allowed. For RQ2, it is compared with ROSE~\cite{zimmermann2004rose}, run with its original
settings (rules with support $\ge 1$ and confidence $\ge 0.1$, top 10, the type rule for new
blocks), and with \emph{equal weights}, the average of the same twelve signals, i.e.\ the
unlearned combination that GA-IaCRec uses for its two heuristics.

\paragraph{Statistics.} Following the guidelines of Arcuri and Briand for randomized
algorithms~\cite{arcuri2011stats}, methods are compared on per-repository values ($n = 53$) with
the Wilcoxon signed-rank test, Holm's correction for multiple comparisons~\cite{holm1979}, and the
Vargha--Delaney $\hat{A}_{12}$ effect size~\cite{vargha2000a12}, read as negligible, small, medium
or large at 0.56, 0.64 and 0.71. We also report bootstrap 95\% confidence intervals of the mean.
